\documentclass[a4paper,12pt]{article}

\usepackage[utf8]{inputenc}
\usepackage[brazil]{babel}
\usepackage{multicol}
\usepackage{enumitem}
\usepackage[a4paper,margin=1.5cm]{geometry}
\usepackage{array}
\usepackage{graphicx}
\usepackage{xcolor}
\usepackage{amsmath}
\usepackage{amssymb}

\begin{document}
\pagestyle{empty}

% =========================================================
% CABEÇALHO
% =========================================================

\begin{table}[h]
\centering
\renewcommand{\arraystretch}{0.6}

\begin{tabular}{|>{\centering\arraybackslash}m{0.20\textwidth}|>{\centering\arraybackslash}m{0.73\textwidth}|}
\hline

\includegraphics[width=2.8cm]{fpf.png}
&
{
\centering
\vspace{0.2cm}
\textbf{\scriptsize CURSO TÉCNICO EM AUTOMAÇÃO INDUSTRIAL} \\[0.15cm]
\textbf{\small Disciplina: Eletrônica Digital 2} \\[0.15cm]
\textbf{\small AVALIAÇÃO DE RECUPERAÇÃO}
}

\tabularnewline
\hline
\end{tabular}

\end{table}

\vspace{0.1cm}

% =========================================================
% DADOS DO ALUNO
% =========================================================

\begin{table}[h]
\centering
\renewcommand{\arraystretch}{1.5}

\begin{tabular}{|>{\raggedright\arraybackslash}m{0.70\textwidth}|>{\raggedright\arraybackslash}m{0.23\textwidth}|}
\hline

\textbf{Aluno(a):} & \textbf{Data:} \\
\hline

\textbf{Professor(a):} & \textbf{Nota:} \\
\hline

\end{tabular}

\end{table}

\vspace{0.05cm}

% =========================================================
% INSTRUÇÕES
% =========================================================

\noindent
\begin{minipage}{\textwidth}
\scriptsize
\begin{itemize}[itemsep=0pt, topsep=0pt, parsep=0pt]
    \item Responda a prova com caneta azul ou preta e não use corretivo, do contrário o(a) aluno(a) não terá direito a questionamentos sobre a correção da prova.
    \item Marque a alternativa escolhida com preenchimento no Cartão-resposta.
    \item Proibido, durante a prova, o uso de telefone celular e/ou aparelhos eletroeletrônicos, empréstimos de materiais, conversas paralelas, nenhum tipo de consulta e sair da sala de aula sem autorização.
    \item O(a) aluno(a) receberá nota 0,0 (Zero), caso seja surpreendido em comunicação com outro aluno de forma verbal ou escrita.
\end{itemize}
\end{minipage}

\vspace{0.5cm}
\noindent\rule{\textwidth}{1pt}

% =========================================================
% PROVA
% =========================================================

\begin{multicols}{2}
\raggedcolumns

\textbf{Cartão-resposta:}

\vspace{0.2cm}

\begin{tabular}{|c|c|c|c|c|c|}
\hline
\textbf{Questões} & A & B & C & D & E \\
\hline
01 & O & O & O & O & O \\
\hline
02 & O & O & O & O & O \\
\hline
03 & O & O & O & O & O \\
\hline
04 & O & O & O & O & O \\
\hline
05 & O & O & O & O & O \\
\hline
06 & O & O & O & O & O \\
\hline
07 & O & O & O & O & O \\
\hline
08 & O & O & O & O & O \\
\hline
09 & O & O & O & O & O \\
\hline
10 & O & O & O & O & O \\
\hline
\end{tabular}

\vspace{0.7cm}

\begin{enumerate}[leftmargin=*]

% =========================================================
% QUESTÃO 1
% =========================================================

\item \textbf{(1,0 ponto)} Sobre os modos de transferência de dados, assinale a alternativa que descreve corretamente a principal diferença entre transferência paralela e serial:

\begin{enumerate}[label=(\Alph*)]
    \item A transferência paralela envia um bit por vez, enquanto a serial envia múltiplos bits simultaneamente.
    \item A transferência serial é sempre mais lenta que a paralela, independentemente da distância.
    \item A transferência paralela envia múltiplos bits simultaneamente através de vários fios, enquanto a serial envia um bit por vez através de um único canal.
    \item A transferência serial exige mais fios de conexão que a paralela.
    \item Não há diferença prática entre os dois modos de transferência.
\end{enumerate}

\vspace{0.2cm}

% =========================================================
% QUESTÃO 2
% =========================================================

\item \textbf{(1,0 ponto)} Em um sistema digital que sincroniza um sinal de dados de 1 MHz com um clock de 10 MHz, qual é a função principal do armazenamento temporário (buffer)?

\begin{enumerate}[label=(\Alph*)]
    \item Aumentar a frequência do sinal de dados para igualar a do clock.
    \item Armazenar os dados até que estejam sincronizados com o clock do sistema, prevenindo perda de informações.
    \item Converter o sinal analógico em digital.
    \item Eliminar completamente o ruído do sinal.
    \item Reduzir a frequência do clock para economizar energia.
\end{enumerate}

\vspace{0.2cm}

% =========================================================
% QUESTÃO 3
% =========================================================

\item \textbf{(1,0 ponto)} Complete a tabela comparativa entre transferência síncrona e assíncrona. Assinale a alternativa que preenche corretamente as lacunas:

\begin{tabular}{|l|l|l|}
\hline
\textbf{Característica} & \textbf{Síncrona} & \textbf{Assíncrona} \\
\hline
Clock dedicado & Necessário & \\
\hline
Bits de start/stop & & Necessários \\
\hline
\end{tabular}

\begin{enumerate}[label=(\Alph*)]
    \item Não necessário / Não necessário
    \item Não necessário / Não utilizado
    \item Utilizado / Não utilizado
    \item Não utilizado / Utilizado
    \item Opcional / Sempre utilizado
\end{enumerate}

\vspace{0.2cm}

% =========================================================
% QUESTÃO 4
% =========================================================

\item \textbf{(1,0 ponto)} Um registrador de deslocamento SIPO (Serial-In Parallel-Out) de 4 bits recebe a sequência de entrada serial \textbf{1011} (primeiro bit 1, depois 0, depois 1, depois 1), começando com todas as saídas em zero. Qual será o estado final das saídas Q3, Q2, Q1, Q0 após 4 ciclos de clock?

\begin{enumerate}[label=(\Alph*)]
    \item Q3=0, Q2=0, Q1=0, Q0=0
    \item Q3=1, Q2=0, Q1=1, Q0=1
    \item Q3=1, Q2=1, Q1=0, Q0=1
    \item Q3=1, Q2=1, Q1=1, Q0=0
    \item Q3=0, Q2=1, Q1=0, Q0=1
\end{enumerate}

\vspace{0.2cm}

% =========================================================
% QUESTÃO 5
% =========================================================

\item \textbf{(1,0 ponto)} Um registrador de deslocamento 74HC595 de 8 bits é utilizado para controlar 8 LEDs. Quantos ciclos de clock são necessários para carregar um byte completo neste registrador?

\begin{enumerate}[label=(\Alph*)]
    \item 4 ciclos
    \item 8 ciclos
    \item 16 ciclos
    \item 1 ciclo
    \item 256 ciclos
\end{enumerate}

\vspace{0.2cm}

% =========================================================
% QUESTÃO 6
% =========================================================

\item \textbf{(1,0 ponto)} Um multivibrador astável usando o CI 555 tem os seguintes componentes: R1 = 1 kΩ, R2 = 10 kΩ e C = 100 μF. Considerando as fórmulas $t_1 = 0,693 \times (R1 + R2) \times C$ para o tempo HIGH e $t_2 = 0,693 \times R2 \times C$ para o tempo LOW, qual é o período total de oscilação (T = t1 + t2)?

\begin{enumerate}[label=(\Alph*)]
    \item Aproximadamente 0,7 ms
    \item Aproximadamente 7,6 ms
    \item Aproximadamente 8,3 ms
    \item Aproximadamente 10,0 ms
    \item Aproximadamente 15,0 ms
\end{enumerate}

\vspace{0.2cm}

% =========================================================
% QUESTÃO 7
% =========================================================

\item \textbf{(1,0 ponto)} Um circuito de temporização precisa gerar um pulso único de 5 segundos cada vez que um botão é pressionado. Qual tipo de multivibrador deve ser utilizado?

\begin{enumerate}[label=(\Alph*)]
    \item Astável, pois oscila continuamente.
    \item Monostável, pois gera um pulso fixo quando acionado.
    \item Biestável, pois mantém dois estados estáveis.
    \item Astável ou Monostável, ambos funcionam.
    \item Nenhum dos anteriores, é necessário um circuito diferente.
\end{enumerate}

\vspace{0.2cm}

% =========================================================
% QUESTÃO 8
% =========================================================

\item \textbf{(1,0 ponto)} Um conversor ADC de 8 bits tem referência de tensão Vref = 5V. Qual é a resolução deste ADC (tensão por passo)?

\begin{enumerate}[label=(\Alph*)]
    \item 5,0 mV
    \item 9,77 mV
    \item 19,53 mV
    \item 39,06 mV
    \item 100,0 mV
\end{enumerate}

\vspace{0.2cm}

% =========================================================
% QUESTÃO 9
% =========================================================

\item \textbf{(1,0 ponto)} Um conversor DAC de 10 bits tem Vref = 3,3V. Para obter uma tensão de saída de aproximadamente 1,65V (metade da referência), qual valor digital deve ser enviado ao DAC?

\begin{enumerate}[label=(\Alph*)]
    \item 256
    \item 512
    \item 768
    \item 1023
    \item 0
\end{enumerate}

\vspace{0.2cm}

% =========================================================
% QUESTÃO 10
% =========================================================

\item \textbf{(1,0 ponto)} Um circuito Sample and Hold (S/H) é utilizado antes de um conversor ADC em um sistema de aquisição de dados. Qual é a principal função deste circuito?

\begin{enumerate}[label=(\Alph*)]
    \item Amplificar o sinal analógico antes da conversão.
    \item Filtrar o ruído do sinal analógico.
    \item Manter a tensão de entrada constante durante o tempo de conversão do ADC.
    \item Converter o sinal digital de volta para analógico.
    \item Aumentar a frequência de amostragem do sistema.
\end{enumerate}


\end{enumerate}

\end{multicols}

\vfill

\end{document}
